\documentclass[11pt,a4paper]{article}

% ── Layout ──
\usepackage[margin=2.4cm]{geometry}
\usepackage{microtype}
\usepackage{enumitem}
\setlist{nosep,leftmargin=*}

% ── Typography ──
\usepackage[T1]{fontenc}
\usepackage{newtxtext,newtxmath}

% ── Bibliography ──
\usepackage[backend=bibtex,style=ieee,sorting=none,maxbibnames=6]{biblatex}
\addbibresource{references.bib}

% ── Graphics & Colors ──
\usepackage{graphicx}
\usepackage[table]{xcolor}
\usepackage{float}
\usepackage{subcaption}

% ── Mathematics ──
\usepackage{amsmath}

% ── Tables ──
\usepackage{booktabs}
\usepackage{multirow}
\usepackage{array}

% ── Hyperlinks ──
\usepackage{hyperref}
\hypersetup{colorlinks=true,linkcolor=blue!50!black,citecolor=blue!50!black,urlcolor=blue!50!black}

% ── Headers ──
\usepackage{fancyhdr}
\setlength{\headheight}{13pt}
\addtolength{\topmargin}{-1pt}
\pagestyle{fancy}
\fancyhf{}
\fancyhead[L]{\small\itshape When Contrastive Learning Meets a Vision Foundation Model}
\fancyhead[R]{\small\thepage}
\renewcommand{\headrulewidth}{0.3pt}

% ── Captions ──
\usepackage{caption}
\captionsetup{font=small,labelfont={bf,small},skip=4pt}

% ── Section formatting ──
\usepackage{titlesec}
\titleformat{\section}{\Large\bfseries}{\thesection}{1em}{}
\titlespacing{\section}{0pt}{12pt}{4pt}
\titleformat{\subsection}{\large\bfseries}{\thesubsection}{1em}{}
\titlespacing{\subsection}{0pt}{10pt}{3pt}

% ── Compact equations ──
\setlength{\abovedisplayskip}{4pt}
\setlength{\belowdisplayskip}{4pt}
\setlength{\abovedisplayshortskip}{2pt}
\setlength{\belowdisplayshortskip}{2pt}

% ── Custom ──
\newcommand{\dinov}{DINOv3\ }
\newcommand{\infonce}{InfoNCE\ }
\newcommand{\methodname}[1]{\textsc{#1}}

\begin{document}

% ═══════════════════════════════════════════════════════════════════
% TITLE PAGE
% ═══════════════════════════════════════════════════════════════════
\begin{titlepage}
\centering
\vspace*{3.5cm}

{\fontsize{26}{32}\selectfont\bfseries When Contrastive Learning\\[6pt]
Meets a Vision Foundation Model\par}
\vspace{10pt}
{\LARGE\itshape A Doomed Fine-Tuning Experiment\\[3pt]
and Its Profound Lessons\par}

\vspace{28pt}
{\large DINOv3 Local Feature Matching for Camera Pose Estimation\\[4pt]
Final Technical Report\par}

\vspace{32pt}
{\large The Chinese University of Hong Kong, Shenzhen\\[3pt]
Image Processing and Computer Vision --- Course Project\\[4pt]
April 2026\par}

\vspace{28pt}
\begin{minipage}{0.88\textwidth}
\large
We attempted to fine-tune DINOv3---a vision foundation model pre-trained on 170 million images---using contrastive learning to improve its performance on local feature matching and camera pose estimation.
Across three paradigms, two architectures, three loss functions, two datasets, and seven systematic experiments, we found that this approach is \textbf{methodologically infeasible}.
The failure arises from a fundamental contradiction: DINOv3's semantic features and the geometric features required for matching are inherently different entities, and geometric distinctiveness has already been irreversibly eliminated as ``noise'' during pre-training.
This report documents the complete experimental journey, the root-cause diagnosis spanning seven layers of analysis, and the definitive conclusion that proving \emph{why} something cannot work holds equal scientific value to proving \emph{what} can.
\end{minipage}
\vfill
\end{titlepage}

\tableofcontents
\newpage

% ═══════════════════════════════════════════════════════════════════
% 1. INTRODUCTION
% ═══════════════════════════════════════════════════════════════════
\section{Introduction}

\subsection{Background}

Self-supervised vision models such as DINOv2 and DINOv3 have demonstrated remarkable capabilities~\cite{dino,dinov2}: without manual annotation, they learn to understand object boundaries, semantic categories, and even 3D structure within images.

DINOv3~\cite{dinov2,darcet2023registers} is a Vision Transformer (ViT-L/16) released by Meta AI, pre-trained on approximately 170 million images. Its self-supervised training objective yields features with strong \emph{semantic smoothness}: patches of a ``door handle'' resemble other ``door handle'' patches; patches of a ``white wall'' resemble patches of the same white wall from a different viewpoint. This property gives DINOv3 non-trivial zero-shot matching capability.

\subsection{Problem Statement}

The problem we address is:

\begin{quote}
\bfseries Given two photographs of the same scene taken from different viewpoints, can we automatically find pixel-level correspondences between them and recover the relative camera pose?
\end{quote}

This is a cornerstone problem in 3D reconstruction, visual localization, and SLAM. The evaluation pipeline proceeds as follows: (1)~establish dense patch correspondences via mutual nearest neighbor (MNN) matching; (2)~estimate the Essential Matrix $E$ from correspondences using USAC~\cite{usac}; (3)~decompose $E$ into relative rotation $R$ and translation $t$ via SVD and cheirality check~\cite{nister2004}; (4)~compute angular error between estimated and ground-truth pose; (5)~report AUC@$5^\circ$, AUC@$10^\circ$, AUC@$20^\circ$ and \emph{Precision} (fraction of matches satisfying the epipolar constraint within $5\times10^{-4}$).

\subsection{Hypothesis (Later Falsified)}

\begin{quote}
\bfseries DINOv3's zero-shot features already produce usable matches (Precision $\approx$\,32--49\%). If we use geometric correspondences as a supervisory signal to fine-tune the model with contrastive learning, matching accuracy should improve further.
\end{quote}

\subsection{Result Overview}

\begin{center}
\color{red!75!black}\bfseries
Every fine-tuning run degraded performance. The best fine-tuned result remained 10\% below the zero-shot baseline. Seven experiments, seven failures, three paradigms, zero exceptions.
\end{center}

% ═══════════════════════════════ 2. TECHNICAL FOUNDATION ═══════════
\section{Technical Foundation}

\subsection{DINOv3 Architecture}

DINOv3 employs a ViT-L/16 backbone~\cite{vit}. For a $448\times448$ input, it partitions the image into $28\times28 = 784$ patches, each producing a 1024-dimensional feature vector.

\begin{table}[t]
\centering\small
\caption{DINOv3 ViT-L/16 specifications.}
\label{tab:dinov3_params}
\begin{tabular}{@{}ll@{}}
\toprule
\textbf{Property} & \textbf{Value} \\
\midrule
Architecture      & ViT-Large (24 blocks, 16 heads) \\
Parameters        & $\sim$300M \\
Feature dimension & 1024 \\
Patch size        & $16\!\times\!16$ px \\
Pre-training data & $\sim$170M images (LVD-1689M) \\
Pre-training      & DINO + iBOT + Register Tokens \\
Position encoding & 2D RoPE \\
Regularization    & LayerScale \\
\bottomrule
\end{tabular}
\end{table}

Key architectural components: \textbf{2D RoPE} encodes patch $(row, col)$ positions by rotating query and key vectors, supporting native extrapolation. \textbf{LayerScale} scales each sub-layer output by a learnable channel-wise factor initialized to $\varepsilon = 10^{-5}$. \textbf{Register Tokens} (4 learnable tokens) absorb global information, letting patch tokens focus on local structure; discarded during extraction.

\subsection{Self-Attention}
\begin{equation}
\text{Attention}(Q, K, V) = \text{softmax}\!\left(\frac{QK^\top}{\sqrt{d_k}}\right)V
\end{equation}
where $Q = XW^Q$, $K = XW^K$, $V = XW^V$.

\subsection{LoRA: Low-Rank Adaptation}

LoRA~\cite{lora} injects trainable low-rank matrices into frozen weights:
\begin{equation}
h = Wx + BA\,x
\end{equation}
where $W \in \mathbb{R}^{d \times d}$ is frozen, $A \in \mathbb{R}^{r \times d}$ is Kaiming-initialized, $B \in \mathbb{R}^{d \times r}$ is zero-initialized, $r \ll d$ ($r=4$). \textbf{Critical}: $B=0$ ensures $\Delta W = 0$ at initialization---the model initially produces \emph{exactly} the zero-shot output. LoRA was injected into Q, K, V, and output projection layers across all 24 blocks, yielding $\sim$0.39M trainable parameters (0.13\% of the backbone).

\subsection{Contrastive Learning and InfoNCE}

Contrastive learning~\cite{cpc,simclr} pulls positive pairs together and pushes negatives apart. \infonce is the dominant formulation:
\begin{equation}
\mathcal{L}_i = -\log\frac{\exp(s_i^+ / \tau)}{\exp(s_i^+ / \tau) + \sum_{j=1}^{K} \exp(s_j^- / \tau)}
\end{equation}
where $s_i^+$ is the positive-pair cosine similarity, $s_j^-$ are $K$ negative similarities, and $\tau$ is the temperature.

\textbf{The $\tau = 0.07$ problem}: Similarities are amplified by $1/0.07 \approx 14.3\times$, causing softmax to approximate argmax---only the single most confusable negative receives meaningful gradient. When the feature space fully collapses to random noise, \infonce converges to its theoretical expectation:
\begin{equation}
\mathbb{E}[\mathcal{L}] = \ln(K+1)
\end{equation}
With $K = 128$ negatives: $\ln(129) \approx 4.8598$. Observing this exact value constitutes \textbf{mathematical proof of mode collapse}.

\subsection{Datasets}

\begin{table}[t]
\centering\small
\caption{Dataset characteristics.}
\label{tab:datasets}
\begin{tabular}{@{}lll@{}}
\toprule
\textbf{Property} & \textbf{NAVI~\cite{navi}} & \textbf{ScanNet~\cite{scannet}} \\
\midrule
Scene type      & Single-object multi-view & Indoor scene scans \\
Test pairs      & 3000 & 1500 \\
Training pairs  & 5596 & 1500 (reuses test pairs) \\
Resolution      & Max dim 1024 & $640\!\times\!480$ \\
Depth maps      & Available (precise 3D) & Unavailable (epipolar) \\
\bottomrule
\end{tabular}
\end{table}

% ═══════════════════════════════ 3. EXPERIMENTAL JOURNEY ═══════════
\section{Experimental Journey}

We report all seven experiments chronologically. Each is described with its hypothesis, method, result, and diagnosis.

\subsection{Phase 0: Zero-Shot Baseline}

\textbf{Method.} Extract 1024-dim features from frozen DINOv3; MNN matching; USAC + 5-point pose estimation.

\textbf{Results.} NAVI: AUC@5\,$=$\,0.24, AUC@10\,$=$\,1.07, AUC@20\,$=$\,3.31, \textbf{Precision\,$=$\,49.17\%}; ScanNet: AUC@5\,$=$\,0.23, AUC@10\,$=$\,1.18, AUC@20\,$=$\,4.44, \textbf{Precision\,$=$\,32.20\%}.

\textbf{Analysis.} Semantic smoothness enables non-trivial matching without task-specific training. This is our baseline.

\subsection{Phase 1: Projection Head + InfoNCE --- Catastrophic Collapse}

\textbf{Hypothesis.} A learnable projection head (1024$\rightarrow$640$\rightarrow$256) trained with \infonce can map DINO features into a more discriminative matching space.

\textbf{Method.} Freeze backbone; add randomly initialized projection head; train with \infonce ($\tau=0.07$, $K=128$).

\textbf{Results.} NAVI Precision: 49.17\% $\rightarrow$ \textbf{3.16\%} ($\downarrow$93.6\%); ScanNet Precision: 32.20\% $\rightarrow$ \textbf{3.47\%} ($\downarrow$89.2\%).

\textbf{Diagnosis.} The training loss converged exactly to $\ln(129) \approx 4.8598$---the theoretical random-guessing limit. This is \textbf{mathematical proof of mode collapse}: all features collapsed to a single point on the hypersphere.

\begin{figure}[H]
\centering
\includegraphics[width=0.78\textwidth]{figures/loss_convergence.png}
\caption{Left: Projection Head loss converges to $\ln(129)$, the signature of mode collapse. Right: Safe Radius keeps loss above the collapse threshold.}
\label{fig:loss}
\end{figure}

\textbf{Lesson.} Randomly initialized components destroy pre-trained features. Fine-tuning must guarantee the starting point equals zero-shot quality.

\subsection{Phase 2: LoRA + InfoNCE --- Zero-Init Is Not Enough}

\textbf{Hypothesis.} LoRA with $B=0$ initialization guarantees training starts from zero-shot quality; it should only improve or maintain performance.

\textbf{Method.} LoRA adapters ($r=4$) on all QKV layers; train with \infonce ($\tau=0.07$).

\textbf{Results.} ScanNet Precision: 32.20\% $\rightarrow$ \textbf{5.51\%} ($\downarrow$82.9\%). Model collapsed despite zero-init guarantee.

\textbf{Diagnosis.} \infonce itself is incompatible with DINOv3's feature space. It demands that \emph{all} non-corresponding patches be pushed apart, but DINOv3's space naturally places semantically similar patches close together. The model destroys its own semantic structure to satisfy the loss.

\subsection{Phase 3: Safe Radius --- The Closest Brush with Success}

\textbf{Hypothesis.} Spatially adjacent patches lie on the same surface and \emph{should} have similar features. Excluding them from the negative set prevents destructive repulsion.

\textbf{Method.} Introduce a \textbf{Safe Radius} of 5 patches: negatives within spatial distance 5 from the anchor are masked out of the contrastive computation.

\textbf{Results.} ScanNet Precision recovered from 5.51\% to \textbf{28.83\%}---only 10.5\% below the zero-shot 32.20\%. This is the single best fine-tuned result.

\textbf{Analysis.} Safe Radius proved ``unreasonable negatives'' were a major degradation cause. Two limitations: (1)~only spatial proximity is addressed, not semantic similarity (distant white-wall patches still repelled); (2)~Precision recovery did not translate to AUC recovery (AUC@20\,$=$\,1.84 vs.\ zero-shot 4.44).

\subsection{Phase 4: StableMatchingLoss --- Removing the Softmax}

\textbf{Hypothesis.} Softmax competition in \infonce is the root problem. A margin-based loss should be gentler.

\textbf{Method.} Moved to RTX 5090 (32\,GB, batch\_size\,$=$\,8). Designed \methodname{StableMatchingLoss}:
\begin{equation}
\mathcal{L} = \underbrace{(1-s_{\text{pos}})^2}_{\text{attraction}} +\; \lambda_{\text{neg}}\underbrace{\text{ReLU}(m - s_{\text{pos}} + s_{\text{hardest}})}_{\text{hardest negative only}} +\; \lambda_{\text{div}}\mathcal{L}_{\text{div}} + \lambda_{\text{reg}}\underbrace{\|\mathbf{f}_{\text{ft}} - \mathbf{f}_{\text{zs}}\|^2}_{\text{feature preservation}}
\end{equation}

\textbf{Results.} Worse than LoRA + Safe Radius: NAVI Precision 5.13\%, ScanNet 10.97\%.

\textbf{Lesson.} The problem is not the loss function's form, but the optimization objective itself. Improving cosine similarity bears no differentiable relationship to improving pose accuracy---the MNN, USAC, and decomposition steps are invisible during training.

\begin{center}
\fbox{\begin{minipage}{0.85\textwidth}
\centering\bfseries\small
All engineering variables eliminated.\\
The problem must be methodological.
\end{minipage}}
\end{center}

\subsection{Phase 5: Route A --- Matchability Predictor}

\textbf{Paradigm shift.} All previous methods modified DINO's feature space. Route A keeps features frozen and trains only a tiny network ($\sim$0.28M params):
\begin{equation}
\text{score}_i = \sigma\bigl(\text{MLP}(\mathbf{f}_i^{\text{DINO}})\bigr) \in [0, 1]
\end{equation}
During inference, only high-scoring patches participate in MNN matching.

\textbf{Results.} NAVI Precision 5.55\%, ScanNet 5.95\%---the predictor filtered out \emph{good} patches rather than bad ones.

\textbf{Failure root cause.} A single patch's feature vector does not encode how \emph{unique} it is within the image. Two white-wall patches have nearly identical features; the predictor cannot distinguish them. This directly motivated Route B.

\subsection{Phase 6: Route B --- Attention Matching Head}

\textbf{Method.} A SuperGlue-style~\cite{superglue} architecture with frozen DINO features: (1)~DINOv3 extracts 1024-dim features; (2)~linear projection $1024\rightarrow256$; (3)~add 2D sinusoidal positional encoding; (4)~4 alternating blocks of Self-Attention $\rightarrow$ Cross-Attention $\rightarrow$ FFN; (5)~L2 normalization; (6)~\textbf{Dual-Softmax Matching}.

Total parameters: $\sim$1M. Training objective: Dual-Softmax NLL Loss. On $784\!\times\!784$ score matrices, the random baseline is $2\times\ln(784) \approx 13.33$.

\textbf{Results: All metrics exactly 0.00\%.}
NAVI: Training loss 13.33 $\rightarrow$ 12.15 (8.9\% reduction; correct-match probability from $1.6\!\times\!10^{-6}$ to $5.2\!\times\!10^{-6}$---still random under argmax).
ScanNet: Training loss 13.33 $\rightarrow$ 13.33 (\textbf{zero learning}). This is the most complete failure.

\textbf{Three-layer diagnosis:}
\begin{enumerate}
    \item \textbf{Extreme label sparsity.} Of $784^2 = 614{,}656$ possible pairs, only $\sim$150 have labels (0.02\%). The remaining 614{,}500 positions drift unconstrained.
    \item \textbf{Zero geometric distinctiveness.} Fifty white-wall patches occupy nearly identical positions in semantic space. Their attention scores are nearly equal, producing uniform attention distributions with zero discriminative power.
    \item \textbf{Data starvation.} 1M-parameter attention model trained on 5,596 pairs (NAVI) and 1,500 pairs (ScanNet). ScanNet's frozen loss (13.33$\rightarrow$13.33) directly proves insufficient data.
\end{enumerate}

% ═══════════════════════════════ 4. COMPLETE RESULTS ═══════════════
\section{Complete Results}

\begin{table}[H]
\centering\small
\caption{All seven experiments: raw results.}
\label{tab:all_results}
\begin{tabular}{@{}cllcccc@{}}
\toprule
\# & \textbf{Method} & \textbf{Dataset} & \textbf{AUC@5} & \textbf{AUC@10} & \textbf{AUC@20} & \textbf{Precision} \\
\midrule
\multirow{2}{*}{0} & \multirow{2}{*}{\color{green!55!black}Zero-Shot (baseline)}
    & NAVI    & 0.24 & 1.07 & 3.31 & \color{green!55!black}49.17\% \\
  & & ScanNet & 0.23 & 1.18 & 4.44 & \color{green!55!black}32.20\% \\
\midrule
\multirow{2}{*}{1} & \multirow{2}{*}{Projection Head + InfoNCE}
    & NAVI    & 0.00 & 0.03 & 0.10 & 3.16\% \\
  & & ScanNet & 0.00 & 0.07 & 0.25 & 3.47\% \\
\midrule
2 & LoRA + InfoNCE (no SafeR) & ScanNet & 0.00 & 0.09 & 0.96 & 5.51\% \\
\midrule
3 & \textbf{LoRA + Safe Radius (bs=1)} & \textbf{ScanNet} & 0.00 & 0.26 & 1.84 & \color{orange!75!black}\textbf{28.83\%} \\
\midrule
\multirow{2}{*}{4} & \multirow{2}{*}{LoRA + StableMatchingLoss}
    & NAVI    & 0.00 & 0.02 & 0.11 &  5.13\% \\
  & & ScanNet & 0.04 & 0.26 & 1.43 & 10.97\% \\
\midrule
\multirow{2}{*}{5} & \multirow{2}{*}{Matchability Predictor (A)}
    & NAVI    & 0.00 & 0.02 & 0.17 &  5.55\% \\
  & & ScanNet & 0.00 & 0.00 & 0.20 &  5.95\% \\
\midrule
\multirow{2}{*}{6} & \multirow{2}{*}{\color{red!75!black}Matching Head (B)}
    & NAVI    & \color{red!75!black}0.00 & \color{red!75!black}0.00 & \color{red!75!black}0.00 & \color{red!75!black}0.00\% \\
  & & ScanNet & \color{red!75!black}0.00 & \color{red!75!black}0.00 & \color{red!75!black}0.00 & \color{red!75!black}0.00\% \\
\bottomrule
\end{tabular}
\end{table}

\begin{figure}[H]
\centering
\includegraphics[width=\textwidth]{figures/precision_all_methods.png}
\caption{Matching Precision across all seven methods. Zero-shot baseline (green) is never surpassed. LoRA + Safe Radius (purple) is the single result approaching baseline.}
\label{fig:precision}
\end{figure}

\begin{figure}[H]
\centering
\includegraphics[width=\textwidth]{figures/auc20_all_methods.png}
\caption{Pose AUC@20 across all methods. Every fine-tuning run degrades pose estimation quality.}
\label{fig:auc20}
\end{figure}

\begin{figure}[H]
\centering
\includegraphics[width=\textwidth]{figures/navi_chronological.png}
\caption{NAVI dataset: Precision and AUC@20 in chronological order. Each fine-tuning attempt produces a cliff-like drop.}
\label{fig:navi}
\end{figure}

\begin{figure}[H]
\centering
\includegraphics[width=\textwidth]{figures/scannet_chronological.png}
\caption{ScanNet dataset: chronological trajectory. LoRA + Safe Radius is the only upward deflection.}
\label{fig:scan}
\end{figure}

\begin{figure}[H]
\centering
\includegraphics[width=\textwidth]{figures/degradation_severity.png}
\caption{Precision degradation from zero-shot (percentage points). More sophisticated methods produce more severe degradation.}
\label{fig:degradation}
\end{figure}

\begin{figure}[H]
\centering
\includegraphics[width=0.85\textwidth]{figures/three_paradigms.png}
\caption{Summary across all three paradigms. No approach improved upon zero-shot.}
\label{fig:paradigms}
\end{figure}

% ═══════════════════════════ 5. ROOT CAUSE ANALYSIS ═══════════════
\section{Root Cause Analysis}

Seven systematic experiments converge on one conclusion: the failure is not engineering but fundamental. This section peels back seven layers of causation.

\subsection{Layer 1: Semantic Features $\neq$ Geometric Features}

DINOv3 learns \textbf{semantic features}: same-category objects are similar, different-category objects are dissimilar. Matching requires \textbf{geometric features}: each specific 3D point must be \emph{unique enough} to distinguish it from all others, even within the same semantic category.

\begin{table}[t]
\centering\small
\caption{Semantic vs.\ geometric features.}
\label{tab:sem_vs_geo}
\begin{tabular}{@{}lll@{}}
\toprule
\textbf{Property} & \textbf{Semantic (DINOv3)} & \textbf{Geometric (matching)} \\
\midrule
Similarity      & ``doorknob $\approx$ doorknob'' & ``\emph{this} corner $\neq$ \emph{that}'' \\
Distribution    & Smooth semantic regions          & Sparse, distinct keypoints \\
Cross-instance  & Generalize across instances      & Discriminate within instance \\
Viewpoint       & Viewpoint-invariant              & Viewpoint-covariant \\
Ideal space     & Continuous manifold              & Discrete clusters \\
\bottomrule
\end{tabular}
\end{table}

\begin{figure}[H]
\centering
\includegraphics[width=0.85\textwidth]{figures/semantic_vs_geometric.png}
\caption{Left: DINOv3's semantic space---clustered by category, smooth. Right: Geometric space---sparse, distinctive, view-covariant.}
\label{fig:semgeo}
\end{figure}

\subsection{Layer 2: Contrastive Learning Destroys Semantic Smoothness}

DINOv3's zero-shot matching works \emph{because of} semantic smoothness: the same object from different views lands nearby in feature space, where nearest-neighbor search finds correct correspondences. Contrastive fine-tuning destroys this: ``pull positives'' forces specific points to jump together; ``push negatives'' tears adjacent points apart. The result: a continuous manifold shredded into discrete clusters. Cross-view generalization is lost.

\textbf{Paradox: the capability you are trying to improve depends on the very property you are destroying.}

\subsection{Layer 3: ``Pulling Positives Together'' Is Itself Wrong}

In dense matching, a ``positive pair''---two patches projecting to the same 3D point from different views---\textbf{should not have identical features}. Illumination differs across views; occlusion states differ; perspective distortion means the $16\!\times\!16$ patch covers slightly different surface areas; attention context differs across images. DINO features naturally encode these viewpoint differences---and this is a \emph{strength}, because viewpoint difference is a critical signal for pose estimation. \textbf{``Pulling positives together'' trains the model to discard the very signal needed for pose estimation.}

\subsection{Layer 4: Three Non-Differentiable Gaps}

\begin{equation}
\text{Cosine sim.} \xrightarrow{\text{non-diff.}} \text{MNN} \xrightarrow{\text{non-diff.}} \text{USAC + E} \xrightarrow{\text{non-diff.}} \text{Pose error}
\end{equation}
The model never sees how improving cosine similarity affects the final pose. ``Better features'' from the training perspective can easily produce ``worse poses'' from the evaluation perspective.

\subsection{Layer 5: The Data Chasm}

\begin{table}[t]
\centering\small
\caption{Pre-training vs.\ fine-tuning data scale.}
\label{tab:data_gap}
\begin{tabular}{@{}lrr@{}}
\toprule
\textbf{Training stage} & \textbf{Data volume} & \textbf{Fraction} \\
\midrule
DINOv3 pre-training & $\sim$170M images & 100\% \\
Fine-tuning (NAVI)  & 5,596 pairs      & $\sim$0.0003\% \\
Fine-tuning (ScanNet) & 1,500 pairs    & $\sim$0.0001\% \\
\bottomrule
\end{tabular}
\end{table}
Every gradient step overwrites vast pre-trained knowledge with extremely limited information. This is not fine-tuning; it is \emph{micro-vandalism}.

\subsection{Layer 6: The Information-Theoretic Ceiling}

This is the deepest finding, crystallized from seven experiments:

\begin{quote}
\bfseries
DINOv3's training objective explicitly minimizes the mutual information between viewpoint variation and semantic features.\\[2pt]
$I(\text{geometric distinctiveness}; \text{semantic feature}) \approx 0$\\[2pt]
Geometric distinctiveness was eliminated as ``noise'' during pre-training. This is not a signal awaiting a better extraction method---it is information that has been \textbf{irreversibly deleted}.
\end{quote}

DINO's self-supervised learning requires ``same object from different viewpoints $\rightarrow$ same feature.'' Viewpoint information---including geometric distinctiveness---is explicitly suppressed. Zero-shot matching achieves Precision 49\%/32\% only because this suppression is imperfect: residual ``edge traces'' of geometric signal survive. Those traces constitute the \textbf{performance ceiling}.

\begin{figure}[H]
\centering
\includegraphics[width=0.78\textwidth]{figures/information_bottleneck.png}
\caption{Information-theoretic view. Geometric information (red) was eliminated during pre-training. No fine-tuning can recover permanently deleted information.}
\label{fig:info}
\end{figure}

\textbf{Analogy.} Consider an aggressively JPEG-compressed image. Post-processing (sharpening, denoising) can improve perceived quality but cannot recover quantized high-frequency details. DINOv3 features are the compressed image; geometric distinctiveness is the lost high-frequency detail; zero-shot matching uses the surviving low-frequency signal. All fine-tuning attempts are efforts to reconstruct high-frequency detail from low-frequency signal---information-theoretically impossible.

\subsection{Layer 7: Sophistication $\propto$ Degradation}

A deeply counter-intuitive finding: \textbf{more sophisticated methods produce more complete degradation}. Safe Radius (a 5-line mask) $\rightarrow$ 28.83\% Precision. StableMatchingLoss (four components) $\rightarrow$ 10.97\%. Matching Head (Self/Cross-Attn + Dual-Softmax) $\rightarrow$ 0.00\%.

This is not coincidence. Sophisticated methods make stronger assumptions about DINO's feature space (e.g., ``attention can distinguish semantically similar but geometrically distinct patches''). Under $I(\text{geometry} \mid \text{semantics}) \approx 0$, every such assumption is false. The sophisticated design becomes a precision destruction mechanism.

% ═══════════════════════════════ 6. CONCLUSION ═══════════════════
\section{Conclusion}

\subsection{Feasibility Verdict}

\begin{quote}
\bfseries\large
Fine-tuning DINOv3's backbone with contrastive learning to improve feature-matching-based camera pose estimation is \textbf{methodologically infeasible}.
\end{quote}

This is not ``we haven't found the right method yet.'' It is ``the right method is to not do this.''

\subsection{Why Zero-Shot Is the Ceiling}

DINOv3's zero-shot features work for matching (Precision 32--49\%) precisely because of their \textbf{semantic smoothness}. Every form of contrastive fine-tuning destroys this smoothness. The loss function demands that the continuous semantic manifold be torn into discrete correspondence clusters. Once smoothness is destroyed, cross-view generalization collapses. More fundamentally: DINOv3's pre-training irreversibly eliminated geometric distinctiveness as noise. Fine-tuning can only operate on information present in the feature space---and that information no longer exists.

\subsection{What Would Work Instead}

\begin{enumerate}
    \item \textbf{Better post-processing on frozen features.} SNN ratio test~\cite{superglue}, improved keypoint selection, RANSAC parameter tuning---zero fine-tuning, direct zero-shot boost.
    \item \textbf{End-to-end differentiable matching with sufficient data.} LoFTR~\cite{loftr} and RoMa~\cite{roma} optimize geometric accuracy directly but require orders of magnitude more training pairs.
    \item \textbf{Purpose-built geometric feature extractors.} SuperPoint~\cite{superpoint}, DISK, ALIKE are designed specifically for matching, unlike general-purpose semantic models.
\end{enumerate}

\subsection{Academic Significance}

\textbf{Understanding why something fails is often more valuable than demonstrating success:} (1)~A success tells you one path works; a failure analysis tells you why hundreds do not. (2)~Our failure is not random but \emph{necessary}---we eliminated every engineering variable and converged on a methodological root cause. (3)~This conclusion guides future researchers: knowing what \emph{not} to do is as valuable as knowing what to do.

\subsection{Methodological Contributions}

This project demonstrates a complete scientific workflow: baseline establishment, hypothesis-driven experimentation, mathematical diagnosis (proving mode collapse via $\ln(129)$), independent algorithmic innovation (Safe Radius), systematic ablation of all engineering variables, and deep methodological reflection.

\begin{center}
\rule{0.45\textwidth}{0.4pt}\\[12pt]
{\large\bfseries Seven experiments, seven layers of understanding,\\[2pt]
converging on a first-principles insight.}\\[10pt]
{\large Proving the infeasibility of a direction holds\\[2pt]
equal academic value to finding a viable one.}\\[10pt]
{\large\color{red!70!black}\bfseries``What \emph{not} to do'' is often more important than ``what to do.''}\\[12pt]
\rule{0.45\textwidth}{0.4pt}
\end{center}

% ═══════════════════════════════ APPENDIX ═══════════════════════
\appendix
\section{Complete Results Table}

\begin{figure}[H]
\centering
\includegraphics[width=\textwidth]{figures/complete_results_table.png}
\caption{All experimental results. Green = zero-shot baseline, red = complete degradation (0.00\%), yellow = partial preservation (28.83\%).}
\label{fig:complete_table}
\end{figure}

% ═══════════════════════════════ REFERENCES ═══════════════════════
\newpage
\printbibliography[heading=bibintoc,title={References}]

\end{document}